\HMTNarrativeHeading{Acción, área e información}

El parámetro de Barbero--Immirzi se presenta sobre los antecedentes
de acción, orientación e incidencia ya construidos. Las coordenadas
angulares y los canales conjugados actúan junto con la inclusión
marcada entre configuraciones hexádicas y octádicas. El funcional
recibe esos datos y determina su evaluación. La relación entre
aritmética y geometría permanece visible en cada uno de sus argumentos.

La estructura de incidencia organiza qué elementos se distinguen y
cómo se prolongan. Una hexada es un soporte de seis posiciones del
diseño; su elevación utiliza una dodecada y un alineamiento declarados.
La octada resultante pertenece a la realización binaria. Los lectores
de área conservan la marca que hace significativa esa correspondencia.
Las figuras ayudan a seguir el transporte de los conjuntos y sus
intersecciones.

La acción y el reloj permiten después expresar la energía asociada
a una lectura. La constante de Boltzmann interviene en la transducción
térmica; el parámetro de Barbero y la unidad areal intervienen en el
incremento geométrico. La constante gravitatoria entra en la
realización que relaciona acción y área de Planck. Al componer estas
expresiones, una misma sección de acción puede cancelarse en el
cociente entre energía e incremento de área. La cancelación exhibe
una dependencia estructural concreta.

La ley de cambio de estado conserva también el término de entropía
relativa y la referencia respecto de la cual se realiza la comparación.
La información de una observación reducida depende de las correlaciones
que mantiene el estado conjunto. Esta distinción prolonga el argumento
del archivo: el contenido de la totalidad y el de un lector parcial
se relacionan mediante operaciones explícitas.

La teoría de conservación adquiere así una aplicación física
determinada. El estado, su orientación, el área que se le asigna y
la lectura térmica forman una composición con antecedentes
identificados. El desarrollo conserva esos antecedentes y sus
pruebas. La síntesis posterior reunirá esta composición con los
balances operatorios y con el ejemplo de entrelazamiento calculable.

